\section{Background Theory}
Mechanical automation has become an essential aspect of modern manufacturing, where repetitive tasks are accomplished through the integration of efficient mechanisms. Among these tasks, counting and packaging operations require accurate motion control, synchronization, and repeatability. Traditionally, such operations are performed manually or with electronically controlled systems, which may increase labor requirements, system complexity, and cost.
\vspace{12px}

This project focuses on the study and analysis of mechanical mechanisms that can be integrated to perform semi-automatic stick counting and packaging operations. The system utilizes purely mechanical components to achieve controlled and sequential motion without relying on complex electronic sensors or controllers.
\vspace{12px}

The theoretical foundation of the project is based on several core areas of mechanical engineering:
\begin{itemize}
    \item Kinematics of Mechanisms: The study of displacement, velocity, and acceleration of moving machine elements such as linkages, cams, ratchet-pawl systems, and intermittent motion mechanisms.
    \item Dynamics of Mechanisms: Investigation of the forces, torques, and power requirements necessary to drive the mechanical system while ensuring smooth and stable operation.
    \item Machine Element Design: Application of engineering principles for the design of shafts, gears, bearings, springs, and other machine components to achieve reliable performance.
    \item Structural Mechanics and Material Selection: Evaluation of stresses and selection of suitable materials to ensure adequate strength, durability, and operational safety.
    \item Computer-Aided Design and Analysis: Three-dimensional modeling and assembly development using FreeCAD, followed by kinematic and dynamic analysis to validate the proposed mechanisms and their interactions.
\end{itemize}

\section{Problem Statement}
Many small and medium-scale manufacturing industries continue to rely on manual methods for counting and packaging stick-type products. Manual operations are repetitive, time-consuming, and prone to inconsistencies in counting accuracy and packaging quality. Although automated systems are available, they often depend on electronic sensors and complex control systems that increase both manufacturing and maintenance costs.
\vspace{12px}

There been currenlty used the drum type stick counter, which is a purely mechanical device that counts sticks as they pass through a rotating drum equipped with counting mechanisms. While this method is simple and does not require electronic components, it has limitations in terms of speed, accuracy, and adaptability to different stick sizes and packaging requirements. The need for a more efficient and versatile solution has led to the exploration of various mechanical mechanisms that can be integrated to achieve semi-automatic stick counting and packaging operations.
\vspace{12px}

From a mechanical engineering perspective, there is a need to study and develop efficient mechanical mechanisms capable of performing these operations through coordinated motion and force transmission. However, the selection, integration, and validation of such mechanisms require detailed kinematic and dynamic investigation.
\vspace{12px}

This project addresses this need by studying and analyzing various mechanical mechanisms, including linkage systems, ratchet-pawl mechanisms, cam mechanisms, and intermittent motion mechanisms, for their application in a semi-automatic stick counting and packaging system. The primary focus is to understand the behavior, motion characteristics, and feasibility of these mechanisms through engineering calculations, CAD modeling, and kinematic-dynamic analysis. 
\vspace{12px}

The ultimate goal is to develop a validated mechanical model that demonstrates the potential for low-cost, efficient automation in stick counting and packaging operations, providing a foundation for future prototype development and further research in this area.